% =====================================================================
%  The Causal Propagation Table
%  Individuation, Contact, and the Accountable Cut in Finite
%  Weighted Graphs
%
%  Self-contained. Every object is a finite weighted graph or a
%  construction on one. No external framework assumed.
% =====================================================================
\documentclass[11pt,twocolumn,a4paper]{article}

\usepackage[utf8]{inputenc}
\usepackage[T1]{fontenc}
\usepackage{lmodern}
\usepackage[a4paper,margin=2.4cm]{geometry}
\usepackage{amsmath,amssymb,amsthm,mathtools}
\usepackage{bm}
\usepackage{mathrsfs}
\usepackage{booktabs}
\usepackage{array}
\usepackage{enumitem}
\usepackage{microtype}
\usepackage{graphicx}
\usepackage{xcolor}
\usepackage[round,sort&compress]{natbib}
\usepackage[colorlinks=true,linkcolor=blue!45!black,
            citecolor=blue!45!black,urlcolor=blue!45!black]{hyperref}
\usepackage{cleveref}

% ---- Theorem environments ----
\newtheoremstyle{thmstyle}{\topsep}{\topsep}{\itshape}{}{\bfseries}{.}{.5em}{}
\theoremstyle{thmstyle}
\newtheorem{theorem}{Theorem}[section]
\newtheorem{lemma}[theorem]{Lemma}
\newtheorem{proposition}[theorem]{Proposition}
\newtheorem{corollary}[theorem]{Corollary}
\theoremstyle{definition}
\newtheorem{definition}[theorem]{Definition}
\newtheorem{axiom}{Axiom}
\newtheorem{construction}[theorem]{Construction}
\newtheorem{example}[theorem]{Example}
\theoremstyle{remark}
\newtheorem{remark}[theorem]{Remark}
\newtheorem{notation}[theorem]{Notation}

\crefname{axiom}{Axiom}{Axioms}
\Crefname{axiom}{Axiom}{Axioms}

% ---- Shorthand ----
\newcommand{\R}{\mathbb{R}}
\newcommand{\N}{\mathbb{N}}
\newcommand{\Z}{\mathbb{Z}}
\newcommand{\Gph}{\Gamma}
\newcommand{\V}{V}
\newcommand{\E}{E}
\newcommand{\wt}{w}
\newcommand{\med}{\mathsf{m}}            % medium vertex
\newcommand{\flo}{\beta}                 % floor
\newcommand{\cut}{\partial}              % edge boundary / cut
\newcommand{\co}{\complement}
\newcommand{\Walk}{W}
\newcommand{\rec}{M}                     % committed record / count
\newcommand{\res}{\rho}                  % residue
\newcommand{\str}{\sigma}                % separation cost
\newcommand{\acc}{\mathcal{A}}           % accountable set
\newcommand{\Rep}{\mathsf{R}}            % representation
\newcommand{\Proc}{\Pi}                  % process graph
\newcommand{\abs}[1]{\lvert#1\rvert}
\newcommand{\norm}[1]{\lVert#1\rVert}
\newcommand{\dd}{\,\mathrm{d}}
\newcommand{\eps}{\varepsilon}

\title{\textbf{The Causal Propagation Table:\\[3pt]
Individuation, Contact, and the Accountable Cut\\
in Finite Weighted Graphs}}

\author{%
Kundai Farai Sachikonye\\
\small Technical University of Munich\\
\small AIMe Registry for Artificial Intelligence\\
\small \texttt{kundai.sachikonye@wzw.tum.de}}

\date{\today}

\begin{document}
\maketitle

% =====================================================================
\begin{abstract}
\noindent
We construct, entirely within the language of finite weighted graphs, a
system in which the act of telling one thing from another---%
\emph{individuation}---is a single, costly, one-use cut, and in which
every physical-seeming notion (identity, time, causation, measurement,
reaction, equilibrium) is a property of cuts and their propagation. The
development rests on one premise, itself derived rather than posited: in a
finite weighted graph whose vertices are the parts individuated within a
whole, every cut separating a part from the rest has weight bounded below
by a strictly positive \emph{floor} $\flo>0$, because identifying a part
is comparing it against a non-completable complement and that comparison
never closes. From the floor we prove, as graph theorems: that a part is
fixed only by its complement (individuation by negation admits no
selector without regress); that the \emph{identity} of a part is not a
vertex label but the invariant minimum-cut value that survives every
relabelling, hence a region and never a point; that the residue of each
cut is the unique edge that makes the next cut possible, so that a
\emph{causal chain} is a walk whose committed-edge count is monotone and
never returns; and that a measurement and a reaction are the \emph{same}
graph operation---a contact edge propagating uncertainty---differing only
in which endpoint is read as instrument, a free labelling of a symmetric
link. We then define the central object, the \emph{causal propagation
table}: a map that takes a designated item vertex, a set of input
vertices, and a process graph, and returns the set of all
\emph{accountable propagations}---walks from input to output whose
composed cut-residue converges, in a bounded $[0,1]$ alignment functional,
to the observed output. We prove the system's three structural freedoms:
\textbf{representation mobility} (the same item admits equivalent
decompositions related by an averaging constraint, so a propagation may
change representation between steps), \textbf{path opacity} (propagations
sharing an endpoint are indistinguishable by any invariant computed at the
endpoint, so the interior path is free), and \textbf{convergence
admissibility} (a propagation is admissible iff it converges to the
output; locally arbitrary, even ``impossible,'' intermediate steps are
permitted, and rejection is only on non-convergence). The output of the
table is shown to be the same kind of object as a cell of the periodic
catalogue of resting items: the catalogue is the table evaluated at the
resting contact, and the table is the catalogue under propagation. We
close by formalising the system's emblem, the single precise cut that can
be made once and dulls the cutter: a cut is a one-use committed edge whose
residue is deposited into the cutter as much as the cut, so no instrument
cuts cleanly and stays sharp---the perfect repeatable traceless cut is
exactly the $\flo=0$ object the floor forbids. Every claim is a theorem
about finite weighted graphs; all interpretation is confined to remarks on
which no theorem depends.
\end{abstract}

\noindent\textbf{Keywords:} finite weighted graph; minimum cut;
individuation by negation; contact floor; causal residue; monotone walk;
path opacity; convergence; representation mobility; accountable
propagation; one-use cut.

\tableofcontents

% =====================================================================
\part{Foundations}
% =====================================================================

% =====================================================================
\section{Introduction}
\label{sec:intro}
% =====================================================================

\subsection{The question and the method}

We ask one question and answer it once: what must hold for a process---a
reaction, a measurement, any sequence of physical happenings---to track an
item from input to output, and what is the output it produces? The answer
we develop is that a process is the propagation of a single repeated
operation, the \emph{cut} that individuates an item from everything else,
and that the output of tracking an item through a process is the same kind
of object as the entry one would write for that item in a catalogue of
resting items.

The method is deliberately austere: every object in the paper is a finite
weighted graph or a construction on one. We use no metric space, no
measure, no probability except as a derived normalisation, and no physical
primitive. The reason is to leave nothing undefined. When we say ``the
identity of an item,'' we mean a specific invariant of a specific graph;
when we say ``a measurement,'' we mean a specific edge; when we say
``causation,'' we mean a specific relation between edges. A reader who
grants only the existence of finite weighted graphs can check every claim.

\subsection{What the paper builds}

The development has four movements.

\emph{Foundations} (Part~I). We fix the contact graph---a finite weighted
graph of individuated parts---and prove that it carries a strictly
positive floor on every cut, that the floor is forced rather than assumed,
that parts are individuated only by negation, and that the identity of a
part is the invariant minimum cut, a region never collapsing to a point.

\emph{Propagation} (Part~II). We prove that the residue of each cut is the
edge that enables the next, so that processes are monotone non-returning
walks; that measurement and reaction are the same symmetric contact
operation; and that the count of committed cuts is the system's time.

\emph{The table} (Part~III). We define the causal propagation table and
prove its three freedoms---representation mobility, path opacity, and
convergence admissibility---culminating in the statement that the
output of the table is a catalogue cell, and that the resting catalogue
(the periodic system) is the table evaluated at rest.

\emph{The cut} (Part~IV). We formalise the system's emblem: the single
precise cut is one-use and self-blunting, and the perfect repeatable cut
is the forbidden $\flo=0$ object. We close with the master equivalence
tying the whole system to the single fact $\flo>0$.

\subsection{The single fact}

Everything below is a consequence of one inequality:
\[
\boxed{\ \flo>0\ }
\]
the floor on the cost of a cut. We first derive it (\Cref{sec:floor}),
then read off the system it forces. The interpretive content---that
$\flo>0$ is why anything is distinct, why processes have a direction, why
no measurement is perfect, and why there is anything to know at all---is
stated in remarks and never used in a proof.

% =====================================================================
\section{The contact graph}
\label{sec:graph}
% =====================================================================

\subsection{Graphs, cuts, and the medium}

\begin{definition}[Finite weighted graph]
\label{def:graph}
A \emph{finite weighted graph} is a triple $\Gph=(\V,\E,\wt)$ with $\V$ a
finite non-empty vertex set, $\E\subseteq\binom{\V}{2}$ an undirected edge
set, and $\wt\colon\E\to\R_{>0}$ a strictly positive weight. We assume
$\Gph$ connected unless stated. The \emph{order} is $\abs{\V}$.
\end{definition}

\begin{definition}[Cut and cut weight]
\label{def:cut}
For $\varnothing\neq S\subsetneq\V$ the \emph{cut} is
$\cut(S)=\{\,e=\{u,v\}\in\E:u\in S,\ v\notin S\,\}$, and its weight is
$\wt(\cut(S))=\sum_{e\in\cut(S)}\wt(e)$. For two disjoint vertex sets
$S,T$ we write $\cut(S,T)$ for the edges with one end in each.
\end{definition}

\begin{definition}[Contact graph and medium]
\label{def:contact-graph}
A \emph{contact graph} is a finite weighted graph $\Gph$ together with a
distinguished vertex $\med\in\V$, the \emph{medium}, adjacent to every
other vertex: $\{v,\med\}\in\E$ for all $v\neq\med$. The vertices
$v\neq\med$ are the \emph{items}; an edge $\{u,v\}$ between two items is a
\emph{contact}, and its weight $\wt(\{u,v\})$ is the cost of the
separation distinguishing $u$ from $v$.
\end{definition}

\begin{remark}[The medium is the rest of the whole, not a background]
\label{rem:medium}
The medium $\med$ is not a spatial container; it is the single vertex
standing for ``everything else,'' against which each item is told apart.
That it is adjacent to every item encodes that every item is, at minimum,
individuated against the whole. No theorem uses any reading of $\med$
beyond \Cref{def:contact-graph}.
\end{remark}

\subsection{Separation cost}

\begin{definition}[Separation cost of an item]
\label{def:sepcost}
The \emph{separation cost} of item $v$ is the minimum weight of a cut
placing $v$ on one side and the medium on the other,
\[
\str(v)=\min_{\substack{S\ni v\\ \med\notin S}}\wt(\cut(S)).
\]
The minimising side $S^{\ast}(v)$ exists and the minimising cut
$\cut(S^{\ast}(v))$ is the \emph{minimum cut} of $v$ against $\med$
(finite $\V$, so the minimum is attained; this is the $s$--$t$ minimum cut
with $s=v$, $t=\med$, computable by max-flow~\citep{fordfulkerson1956,%
mengertheorem}).
\end{definition}

% =====================================================================
\section{The floor, derived}
\label{sec:floor}
% =====================================================================

We show the central constant is not an assumption but a consequence of two
facts: an item is a proper part, and the whole is non-completable.

\subsection{Non-completability}

\begin{definition}[Expanding family; non-completable whole]
\label{def:noncompletable}
An \emph{expanding family} is a sequence of contact graphs
$\Gph_0\subseteq\Gph_1\subseteq\Gph_2\subseteq\cdots$, each obtained from
the last by adding vertices and edges (of weight $\ge\flo$ as fixed below)
while retaining the medium and all prior vertices and edges. The whole is
\emph{non-completable} if the family has no terminal stage: for every
finite $\Gph_k$ there is a $\Gph_{k+1}\supsetneq\Gph_k$. This is an order
condition (no last stage), not a cardinality assumption.
\end{definition}

\begin{definition}[Identification and its residual]
\label{def:identification}
To \emph{identify} item $v$ in $\Gph_k$ is to determine $v$ by its cut
against the medium, $\cut(S^{\ast}(v))$. The identification is
\emph{complete} if every item of the whole has been brought to bear in the
cut---i.e.\ the cut is taken in a graph that contains the whole. The
\emph{identification residual} $\res(v)$ is the weight of the cut content
that a complete identification would remove but a finite-stage
identification cannot; $\res(v)=0$ iff identification is complete.
\end{definition}

\subsection{Floor from non-completability}

\begin{theorem}[Floor from infinitude]
\label{thm:floor-derived}
Let the whole be non-completable and let $v$ be an item (a proper part:
$v\neq\med$, and $v$ does not exhaust $\V$). Then the identification of $v$
cannot be completed at any finite stage, and its residual is strictly
positive: $\res(v)>0$.
\end{theorem}

\begin{proof}
Suppose $\res(v)=0$, i.e.\ identification of $v$ is complete at some
finite stage $\Gph_k$. By \Cref{def:identification} completeness means
every item of the whole has entered the cut of $v$ against $\med$; that
is, the cut is computed in a graph containing the whole. But then
$\Gph_k$ contains the whole, contradicting non-completability
(\Cref{def:noncompletable}), which guarantees a further stage
$\Gph_{k+1}\supsetneq\Gph_k$ with an item not in $\Gph_k$, hence not yet
brought to bear. Therefore no finite stage completes the identification,
and $\res(v)>0$.
\end{proof}

\begin{figure*}[!htbp]
    \centering
    \includegraphics[width=\textwidth]{figures/panel_1_floor.png}
    \caption{\textbf{The Floor: Minimum Cuts over Random Contact Graphs.}
    (\textbf{A})~Three-dimensional scatter of $400$ random contact graphs in (number of items, minimum-cut cost, cut-side size) space, points coloured by cut cost (viridis), with the floor plane at cost $\beta$ drawn in red. Every point lies on or above the floor plane: no item is separated from the medium for less than $\beta$ (Theorem~\ref{thm:floor}).
    (\textbf{B})~Distribution of the exact minimum-cut costs (\texttt{networkx} max-flow) across the $400$ graphs, with the floor $\beta$ marked (red dashed). No mass falls below $\beta$, the absence of any sharp cut (Corollary~\ref{cor:no-sharp}).
    (\textbf{C})~Minimum-cut cost versus number of items (jittered) for all graphs; the entire cloud sits above the floor line $\beta$ (red dashed) independently of graph order, confirming the bound holds at every size.
    (\textbf{D})~Survival curve: the fraction of cuts with cost at least a given threshold, dropping from one only after the threshold passes $\beta$ (red dashed). The fraction at or above $\beta$ is unity---every one of the $300$ validated cuts clears the floor.}
    \label{fig:panel_floor}
\end{figure*}

\begin{theorem}[Positive floor]
\label{thm:floor}
There is a constant $\flo>0$ with $\wt(e)\ge\flo$ for every edge and
$\str(v)\ge\flo$ for every item: no two adjacent vertices are separated at
zero cost, and no item is individuated for free.
\end{theorem}

\begin{proof}
Set $\flo:=\inf_v\res(v)$. By \Cref{thm:floor-derived} each $\res(v)>0$;
the residual is the irreducible boundary content of $v$'s individuation,
which is carried by the edges of $v$'s minimum cut, so each such edge has
weight at least the per-item residual and hence at least $\flo$. As every
edge is the boundary edge of some item's individuation (every edge
$\{u,v\}$ is in $\cut(\{u\})$), every weight is $\ge\flo$. Connectivity
forces $\cut(S^{\ast}(v))\neq\varnothing$ (at least one edge crosses any
$v$--$\med$ separation), so
$\str(v)=\wt(\cut(S^{\ast}(v)))\ge\flo>0$.
\end{proof}

\begin{corollary}[No sharp cut]
\label{cor:no-sharp}
There is no cut of weight $0$ separating a connected contact graph into a
nonempty item set and its nonempty complement. Every individuation
deposits weight $\ge\flo$.
\end{corollary}

\begin{remark}[The floor is the trace of the whole on its parts]
\label{rem:trace}
\Cref{thm:floor-derived,thm:floor} make $\flo$ a derived quantity: it is
the residual the non-completable whole forces on every finite part.
A completable whole would permit $\res(v)=0$, hence $\flo=0$, hence sharp
cuts and point-identities; the impossibility of completing the comparison
is exactly what makes the floor positive. We return to this as the master
equivalence (\Cref{sec:master}). No later theorem assumes $\flo$; each
uses only its positivity.
\end{remark}

% =====================================================================
\section{Individuation by negation}
\label{sec:negation}
% =====================================================================

\subsection{No privileged vertex}

\begin{definition}[Selector; selector-free determination]
\label{def:selector}
A \emph{selector} for an item set $U\subseteq\V$ is a rule fixing $U$ by
naming a distinguished vertex of $U$ in advance. A determination of $U$ is
\emph{selector-free} if it names no vertex of $U$. The graph has \emph{no
privileged vertex} if the data presented to a determining rule are
invariant under the weighted automorphism group $\operatorname{Aut}(\Gph)$,
so no vertex is singled out by the data alone.
\end{definition}

\begin{theorem}[Individuation by negation]
\label{thm:negation}
In a contact graph with no privileged vertex, any selector-free
determination of an item set $U$ fixes $U$ as the complement of its
negation set, $U=\V\setminus\co U$ with $\co U:=\V\setminus U$ specified
first; and conversely every admissible $\co U$ determines a unique $U$. No
selector-based rule fixes $U$ without invoking a further selector.
\end{theorem}

\begin{proof}
\emph{Necessity.} To fix $U$ is to decide membership of each vertex. A
positive rule asserting membership of a distinguished vertex of $U$ must
name that vertex; with $\operatorname{Aut}(\Gph)$-invariant data no vertex
is named, so the rule must itself supply one---a selector---which must in
turn be fixed, either given in advance (excluded) or by a prior rule
facing the same dichotomy. The regress halts only at a rule naming no
vertex of $U$: the withholding of $U$ from the whole, i.e.\ specifying
$\co U=\V\setminus U$ and taking $U$ as the remainder. This names only
``the rest,'' supplied by $\V$. \emph{Sufficiency and uniqueness.}
Complementation $U\mapsto\V\setminus U$ is an involution on subsets of the
finite set $\V$, so $\co U$ determines $U$ uniquely and conversely.
\end{proof}

\begin{corollary}[An item is fixed by its non-neighbours]
\label{cor:vertex-negation}
A single item $v$ is determined selector-free by the partition of the rest
into its neighbours $N(v)$ (what $v$ contacts) and its non-neighbours
$\V\setminus N[v]$ (what it neither is nor contacts), with no intrinsic
mark on $v$. Its identity is carried by this negation structure, not by
its label.
\end{corollary}

\begin{proof}
\Cref{thm:negation} with $U=\{v\}$, refined by adjacency: the
selector-free datum fixing $\{v\}$ is $\V\setminus\{v\}$ split into
$N(v)$ and $\V\setminus N[v]$; no vertex of $\{v\}$ is named.
\end{proof}

\begin{figure*}[!htbp]
    \centering
    \includegraphics[width=\textwidth]{figures/panel_2_negation.png}
    \caption{\textbf{Individuation by Negation: the Complement Involution.}
    (\textbf{A})~Three-dimensional scatter of $200$ random subsets of an eight-element whole in $(|U|,\,|\co U|,\,|\co\co U|)$ space, coloured by $|U|$ (plasma). Every point satisfies $|\co\co U|=|U|$: double complementation returns the part, the involution underlying Theorem~\ref{thm:negation}.
    (\textbf{B})~Subset size $|U|$ versus complement size $|\co U|$; all points lie on the conservation line $|U|+|\co U|=8$ (red dashed): the complement is exactly ``the rest,'' fixed once the whole is given.
    (\textbf{C})~Double-complement size $|\co\co U|$ versus $|U|$ on the identity diagonal (grey dashed); the points coincide with the diagonal, so $\co$ is its own inverse and $\co U$ determines $U$ uniquely.
    (\textbf{D})~Rules required to fix a part as a function of nesting level: a positive (selector-based) rule regresses without bound (red), while negation closes at level one (green). The diverging red curve is the selector regress that Theorem~\ref{thm:negation} excludes; the flat green line is complementation, which needs only the whole.}
    \label{fig:panel_negation}
\end{figure*}

% =====================================================================
\section{Identity is the invariant cut}
\label{sec:identity}
% =====================================================================

We now identify what, about an item, is invariant across all re-encodings
of the graph---and show it is never a vertex label and never a point.

\subsection{The intrinsic minimum cut}

\begin{definition}[Weighted isomorphism; relabelling]
\label{def:iso}
A \emph{weighted isomorphism} $\phi\colon\Gph\to\Gph'$ is a vertex
bijection preserving the medium, edges, and weights:
$\{u,v\}\in\E\iff\{\phi u,\phi v\}\in\E'$ and
$\wt'(\{\phi u,\phi v\})=\wt(\{u,v\})$, with $\phi(\med)=\med'$. A
\emph{relabelling} is an automorphism ($\Gph'=\Gph$).
\end{definition}

\begin{definition}[Intrinsic floor of an item]
\label{def:intrinsic}
For an item $v$ in an expanding family $\Gph_0\subseteq\Gph_1\subseteq
\cdots$, the \emph{intrinsic floor} is
$\flo^{\ast}(v):=\lim_{k\to\infty}\min_{e\ni v}\wt_k(e)$, the limit of the
least weight incident to $v$ as the whole grows.
\end{definition}

\begin{lemma}[Existence and positivity of the intrinsic floor]
\label{lem:intrinsic}
For every item $v$ the limit $\flo^{\ast}(v)$ exists and satisfies
$\flo^{\ast}(v)\ge\flo>0$.
\end{lemma}

\begin{proof}
The sequence $\min_{e\ni v}\wt_k(e)$ is non-increasing (adding edges can
only lower the minimum incident weight) and bounded below by $\flo$
(\Cref{thm:floor}). A non-increasing sequence bounded below converges, to
a limit $\ge\flo$.
\end{proof}

\subsection{Identity is invariant and region-valued}

\begin{theorem}[The minimum cut is the conserved invariant]
\label{thm:invariant}
The separation cost $\str(v)$ and the intrinsic floor $\flo^{\ast}(v)$ are
weighted-graph invariants: if $\phi\colon\Gph\to\Gph'$ is a weighted
isomorphism then $\str_{\Gph'}(\phi v)=\str_{\Gph}(v)$ and
$\flo^{\ast}(\phi v)=\flo^{\ast}(v)$. The vertex label of $v$ is not an
invariant: it is permuted by automorphisms.
\end{theorem}

\begin{proof}
$\phi$ carries the cut $\cut(S^{\ast}(v))$ to a cut of $\phi v$ against
$\med'$ of equal weight (weights preserved), and the bijection on cuts has
inverse $\phi^{-1}$, so the minima agree:
$\str_{\Gph'}(\phi v)=\str_\Gph(v)$. The same applies stagewise in an
expanding family, giving $\flo^{\ast}(\phi v)=\flo^{\ast}(v)$. Labels are
permuted by any nontrivial automorphism, hence not invariant.
\end{proof}

\begin{theorem}[Identity is a region, never a point]
\label{thm:region}
The identity of an item---the invariant data fixing it---is the minimum
cut $\cut(S^{\ast}(v))$, a set of edges of total weight $\str(v)\ge\flo>0$.
It is realised by no single vertex: in general the minimising side
$S^{\ast}(v)$ is not the singleton $\{v\}$, and the identity cannot be
reduced to a zero-content (point) datum.
\end{theorem}

\begin{proof}
By \Cref{thm:invariant} the invariant fixing $v$ is the minimum-cut value
and its cut set, not the (variable) label. By \Cref{thm:floor} the cut has
weight $\ge\flo>0$, so it is never a zero-content point. Non-locality: take
a contact graph with two dense item-clusters joined by a single
$\flo$-weight contact; the minimum cut against $\med$ may separate a whole
cluster (side $S^{\ast}$ of size $>1$) at weight below
$\min_v\wt(\cut(\{v\}))$, so the minimiser is not a singleton and the
identity is borne by a region of edges, not a point.
\end{proof}

\begin{remark}[What a name is]
\label{rem:name}
\Cref{thm:invariant,thm:region} say that a name---a vertex label such as
``this item''---carries no invariant; it is the relabellable surface. What
is invariant is the minimum cut, the region of edges by which the item is
told from the rest, of weight $\ge\flo$. A name denotes a cut at the
current stage of an expanding family; the identity it points at is the
intrinsic floor the cut is converging to and, by \Cref{thm:region}, never
reaches as a point. This is used substantively in \Cref{sec:catalogue}:
a catalogue entry is a name for a cut, not a point.
\end{remark}

\begin{figure*}[!htbp]
    \centering
    \includegraphics[width=\textwidth]{figures/panel_3_identity.png}
    \caption{\textbf{Identity as an Invariant Region, Never a Point.}
    (\textbf{A})~Three-dimensional rendering of the two-cluster contact graph: items (blue if on the source side $S$ of the minimum cut, red otherwise) and the medium (orange) at the origin, lifted on the vertical axis by cut membership, with edges drawn at width proportional to weight. The minimum cut places a multi-item set on the source side---identity is borne by a region, not a single vertex.
    (\textbf{B})~Minimum-cut cost before versus after a random relabelling (weighted isomorphism) of the items, over $120$ graphs; all points lie on the identity diagonal (grey dashed), so the cost is label-independent (Theorem~\ref{thm:invariant}).
    (\textbf{C})~Distribution of minimum-cut side sizes over $400$ random graphs; the mass lies predominantly above one item (the threshold at $1.5$, red dashed), the empirical content of ``the minimising side is in general not a singleton'' (Theorem~\ref{thm:region}).
    (\textbf{D})~Minimum-cut value for four different separated pairs in the two-cluster instance: the cluster-separating pairs (green) realise the cheap inter-cluster cut, while same-cluster pairs (orange) cost more, showing the cut value is a property of the partition, not of a chosen endpoint.}
    \label{fig:panel_identity}
\end{figure*}

% =====================================================================
\part{Propagation}
% =====================================================================

% =====================================================================
\section{The residue is the next cut}
\label{sec:residue}
% =====================================================================

We now make cuts dynamic: a sequence of cuts is a walk, and the residue of
each cut is precisely the edge that the next cut requires.

\subsection{Committed cuts and the record}

\begin{definition}[Committed edge; committed record]
\label{def:committed}
A \emph{cut event} on $\Gph$ commits a contact: it fixes one edge
$e\in\E$ as a realised separation (the boundary just drawn). A
\emph{process} is a finite sequence of cut events
$P=(e_1,e_2,\dots,e_k)$, $e_i\in\E$, possibly with repetition of vertices
but each $e_i$ a distinct committing act. The \emph{committed record}
after $P$ is the count $\rec(P)=k$ of cut events.
\end{definition}

\begin{definition}[Residue of a cut event]
\label{def:residue}
The \emph{residue} of cut event $e_i=\{u,v\}$ is the weight
$\res_i:=\wt(e_i)\ge\flo>0$ it deposits as committed boundary content
between $u$ and $v$. The residue is irreducible (\Cref{thm:floor}): no cut
event deposits less than $\flo$.
\end{definition}

\subsection{Residue as causal propagator}

\begin{theorem}[The residue enables the next cut]
\label{thm:propagator}
In a process $P=(e_1,\dots,e_k)$, the residue $\res_i>0$ of cut $e_i$ is
necessary and sufficient for cut $e_{i+1}$ to be a distinct committing
event: $e_{i+1}$ is a new cut iff the state after $e_i$ differs from the
state before, which holds iff $\res_i>0$. If $\res_i=0$ there is no
difference between pre- and post-$e_i$ states and no ground for a distinct
$e_{i+1}$.
\end{theorem}

\begin{proof}
A cut event $e_{i+1}$ is distinct from its predecessor iff it acts on a
state changed by $e_i$. The only change $e_i$ makes is the deposit of its
residue $\res_i$ (the committed boundary content). If $\res_i>0$ the
post-state carries a committed edge absent from the pre-state, so
$e_{i+1}$ acts on a strictly different graph configuration and is a
distinct event; the residue is the difference enabling it. If $\res_i=0$
(excluded by \Cref{thm:floor} but argued for necessity) the pre- and
post-states coincide and $e_{i+1}$ has no changed state to act on, hence is
not a distinct committing event. Thus $\res_i>0$ is necessary and
sufficient.
\end{proof}

\begin{figure*}[!htbp]
    \centering
    \includegraphics[width=\textwidth]{figures/panel_4_record_residue.png}
    \caption{\textbf{The Monotone Record and the Residue Propagator.}
    (\textbf{A})~Three-dimensional trajectory of a process in (step, committed record $M$, cumulative residue) space: the record climbs monotonically while the residue accumulates, the two coupled because each cut deposits a residue and increments the record (Theorem~\ref{thm:monotone}, Theorem~\ref{thm:propagator}).
    (\textbf{B})~Committed record $M$ versus step for an eight-event process; the staircase strictly increases, and the ``undo'' events (red points) increment the record exactly as forward cuts (blue) do---un-committing is itself a further cut, so the record never returns.
    (\textbf{C})~Whether a distinct next cut is possible (yes/no) versus the residue of the current cut; the transition occurs at residue zero, with points coloured green when the residue clears the floor $\beta$ (red dashed) and red otherwise. Residue $>0$ is necessary and sufficient for the next cut.
    (\textbf{D})~Residue at each step of a five-step chain, bars blue where the residue clears the floor and red where it would vanish; the dashed floor line $\beta$ and the halt marker (grey dotted) show the chain stopping exactly where a residue would fall below $\beta$---no residue, no next event.}
    \label{fig:panel_record_residue}
\end{figure*}

\begin{corollary}[A process is a residue chain]
\label{cor:chain}
A process is a chain in which each cut's residue is the cause of the next:
$\res_1\to e_2,\ \res_2\to e_3,\ \dots$. The chain has a first cut (the
seed) and proceeds only by the deposit of residues; remove any $\res_i$
and the chain halts at step $i$.
\end{corollary}

\begin{proof}
Immediate from \Cref{thm:propagator} applied along $P$.
\end{proof}

\subsection{Monotonicity and non-return}

\begin{theorem}[Monotone non-returning record]
\label{thm:monotone}
The committed record $\rec(P)$ is strictly increasing along any process:
each cut event increments it by one, and no process decreases it. In
particular a process that revisits a vertex configuration does not return
to a prior state: the configuration recurs but the record has strictly
grown.
\end{theorem}

\begin{proof}
By \Cref{def:committed} each cut event is a distinct committing act,
incrementing $\rec$ by one. To decrease $\rec$ a committed edge would have
to be un-committed; un-committing is itself a further committing act (a new
cut redrawing the boundary), which increments $\rec$ rather than
decrementing it---the residue of the original cut cannot be withdrawn
without depositing a new residue. Hence $\rec$ is monotone. Two process
states with the same vertex configuration but different records are
distinguished by the record (the partition of states ``record $\le r$''
versus ``$>r$'' separates them), so they are distinct states: recurrence of
configuration is not return of state.
\end{proof}

\begin{remark}[The record is time]
\label{rem:time}
\Cref{thm:monotone} gives a monotone scalar attached to every process---the
committed count. It is the system's clock: each cut event is one tick, the
count never decreases, and equal configurations at different counts are
different states. ``Before'' and ``after'' are the order on the record.
Nothing in the proof uses any external time parameter; the order is
intrinsic to the sequence of cuts. No theorem below depends on this
reading.
\end{remark}

\subsection{The three readings of a residue}

\begin{theorem}[Time, information, entropy as one residue]
\label{thm:three-readings}
A residue $\res_i$ admits exactly three readings, according to how its link
to the next cut is resolved:
\begin{enumerate}[label=\textup{(\roman*)},leftmargin=2.2em]
\item \textbf{Time}: as the increment of the committed record, $\res_i$ is
one tick (\Cref{thm:monotone}).
\item \textbf{Information}: when $\res_i$ is tracked as the specific cause
of a specific next cut $e_{i+1}$, it is a resolved causal link---one bit of
``which cut caused which''---and the chain is determinate.
\item \textbf{Entropy}: when residues are present but their links to next
cuts are not individually tracked, their aggregate is the Shannon content
$H=-\sum_i p_i\log p_i$ of the unresolved link distribution
\citep{shannon1948,cover2006}.
\end{enumerate}
Tracking a process---resolving which residue causes which next cut---moves
its reading from entropy toward information; losing track moves it back.
\end{theorem}

\begin{proof}
(i) is \Cref{thm:monotone}. (ii): a tracked link is a determinate function
``$\res_i\mapsto e_{i+1}$''; specifying it for each $i$ is specifying the
chain exactly, i.e.\ resolving the causal structure---one bit per binary
link. (iii): with links unresolved, the chain is known only through the
distribution over which residue caused which next cut, whose information
content is $H$ by definition \citep{shannon1948}. Resolving links converts
unknown distribution (entropy) into determinate functions (information);
the conversion is exactly the tracking. \qed
\end{proof}

\begin{remark}[Tracking is the table's work]
\label{rem:tracking}
\Cref{thm:three-readings} is the hinge to Part~III. The causal propagation
table is the device that \emph{tracks}: it resolves causal links along a
process, converting the residue chain from an entropy reading (an opaque
box: things go in, things come out) to an information reading (an accounted
chain: this cut caused that one). The output of the table is the resolved
chain---the residues read as information.
\end{remark}

% =====================================================================
\section{Measurement and reaction are one operation}
\label{sec:contact-symmetry}
% =====================================================================

We prove the identity on which the table rests: a measurement step and a
reaction step are the same graph operation---a contact edge propagating
uncertainty---and ``instrument'' versus ``sample'' is a free labelling of a
symmetric link.

\subsection{The contact operation}

\begin{definition}[Contact operation]
\label{def:contact-op}
A \emph{contact operation} on items $u,v$ is the cut event committing the
edge $\{u,v\}$: it individuates $u$ against $v$ and $v$ against $u$
simultaneously, depositing residue $\wt(\{u,v\})\ge\flo$ on the single
shared boundary. The operation is \emph{symmetric}: it is the same edge
whether read from $u$ or from $v$.
\end{definition}

\begin{theorem}[Instrument--sample symmetry]
\label{thm:symmetry}
A contact operation on $\{u,v\}$ admits two readings---``$u$ measures $v$''
(read $u$ as instrument) and ``$v$ measures $u$'' (read $v$ as
instrument)---which are the same cut event. Neither reading is privileged
by the graph: the edge $\{u,v\}$ is unordered, and exchanging the labels
$u\leftrightarrow v$ is an isomorphism of the operation. Hence
``instrument'' and ``sample'' are a labelling of the two ends of one
symmetric link, not a property of either end.
\end{theorem}

\begin{proof}
The committed object is the undirected edge $\{u,v\}$ with a single weight;
``measure'' in either direction is reading the same residue from one end.
The transposition $u\leftrightarrow v$ fixes the edge and its weight, hence
is an automorphism of the one-edge operation; by \Cref{thm:negation} no
privileged vertex exists to break the symmetry. The two readings therefore
denote the same cut event under a label exchange. \qed
\end{proof}

\begin{theorem}[Reaction = measurement]
\label{thm:reaction-measurement}
Let a \emph{reaction step} be a contact operation committing $\{u,v\}$ that
is not subsequently tracked (its causal link unresolved), and a
\emph{measurement step} be a contact operation committing $\{u,v\}$ that is
tracked (its causal link resolved). Then a reaction step and a measurement
step are the same operation (\Cref{def:contact-op}) differing only in the
reading of its residue (\Cref{thm:three-readings}): reaction reads it as
entropy, measurement as information. There is no graph-level distinction
between them.
\end{theorem}

\begin{proof}
Both are the cut event committing the edge $\{u,v\}$, depositing the same
residue $\wt(\{u,v\})$ (\Cref{def:contact-op}); the operations are
identical as graph acts. The only difference is whether the residue's link
to the next cut is resolved (\Cref{thm:three-readings}): tracked
$\Rightarrow$ information reading $\Rightarrow$ ``measurement''; untracked
$\Rightarrow$ entropy reading $\Rightarrow$ ``reaction.'' Tracking is not a
property of the edge but of the observer's resolution of the chain
(\Cref{rem:tracking}). Hence no graph-level distinction exists. \qed
\end{proof}

\begin{corollary}[Either end is the instrument; synthesis = analysis]
\label{cor:either-end}
For any contact in a process, either endpoint may be read as the
instrument and the other as the sample (\Cref{thm:symmetry}), and the
process may be read forward (building a product) or backward (resolving a
measurement) (\Cref{thm:reaction-measurement}). An ionic contact and a
spectrometric contact are the same operation: a contact edge propagating
uncertainty.
\end{corollary}

\begin{proof}
\Cref{thm:symmetry} gives end-exchange freedom;
\Cref{thm:reaction-measurement} gives reading-direction freedom. Both are
label freedoms on the same symmetric operation. \qed
\end{proof}

\begin{remark}[The mystery box]
\label{rem:box}
A \emph{mystery box} is a process whose contacts are committed but whose
causal links are unresolved---the entropy reading
(\Cref{thm:three-readings}(iii)): inputs go in, an output comes out, and
the interior chain is opaque. By \Cref{thm:reaction-measurement} the
interior is a chain of contact operations indifferent to whether each is
called reaction or measurement. The table's task (Part~III) is to resolve
the interior into an accountable chain---to convert the box from entropy to
information. No theorem depends on the word ``box.''
\end{remark}

% =====================================================================
\part{The Causal Propagation Table}
% =====================================================================

% =====================================================================
\section{Alignment and accountability}
\label{sec:alignment}
% =====================================================================

To say a propagation ``converges to the output'' we need a scalar that
measures how far a graph state is from a target. We build it from cuts, so
it remains internal to the graph language.

\subsection{The alignment functional}

\begin{definition}[Alignment functional]
\label{def:alignment}
Fix a contact graph $\Gph$ and a target item $x^{\ast}$ (the observed
output). The \emph{alignment} of an item $x$ to $x^{\ast}$ is
\[
\str(x,x^{\ast})\ :=\ \min_{\substack{S\ni x\\ x^{\ast}\notin S}}
\wt(\cut(S))\ \in\ [\flo,\,\Omega],
\]
the minimum-cut weight separating $x$ from $x^{\ast}$, where
$\Omega:=\wt(\E)$ is the total weight (a finite upper bound). Normalising,
the \emph{alignment score} is
\[
a(x,x^{\ast})\ :=\ \frac{\str(x,x^{\ast})}{\Omega}\ \in\ (0,1].
\]
$a$ is small when $x$ and $x^{\ast}$ are cheaply separated (nearly the same
cut) and large when they are expensively separated (very different cuts).
\end{definition}

\begin{lemma}[The floor bounds alignment below]
\label{lem:align-floor}
For distinct items $x\neq x^{\ast}$, $a(x,x^{\ast})\ge\flo/\Omega>0$:
no two distinct items are perfectly aligned. Alignment $a=0$ would require
$\str=0$, forbidden by \Cref{thm:floor}.
\end{lemma}

\begin{proof}
Immediate from \Cref{thm:floor}: any $x$--$x^{\ast}$ cut has weight
$\ge\flo$, so $a\ge\flo/\Omega>0$. \qed
\end{proof}

\begin{remark}[Why a positive floor on alignment matters]
\label{rem:align-floor}
\Cref{lem:align-floor} is the convergence analogue of ``identity is never
a point'' (\Cref{thm:region}): a propagation can converge \emph{toward} the
output but never to perfect ($a=0$) alignment, only to within the floor.
``Accountable'' below means ``converged to within the floor,'' never
``identical.''
\end{remark}

\subsection{Accountability}

\begin{definition}[Accountable state]
\label{def:accountable}
A graph state is \emph{accountable to the output $x^{\ast}$ at tolerance
$\eps$} if its current item $x$ satisfies $a(x,x^{\ast})\le\flo/\Omega+\eps$
for the prescribed $\eps\ge0$, i.e.\ it has converged to within the floor
plus tolerance of the output. The least achievable is the floor itself
($\eps=0$): convergence to within the floor.
\end{definition}

% =====================================================================
\section{Processes, propagations, and the table}
\label{sec:table}
% =====================================================================

\subsection{The process graph}

\begin{definition}[Process graph]
\label{def:process-graph}
A \emph{process graph} $\Proc=(\V,\E,\wt;\,I,x^{\ast})$ is a contact graph
together with a set $I\subseteq\V$ of \emph{input items}, and a designated
\emph{output item} $x^{\ast}\in\V$. The edges of $\Proc$ are the available
contact operations (the ``conditions and processes'' the user supplies):
each edge is a permitted cut event.
\end{definition}

\begin{definition}[Propagation]
\label{def:propagation}
A \emph{propagation} of a designated \emph{tracked item} $v_0\in I$ through
$\Proc$ is a walk $W=(v_0,e_1,v_1,e_2,\dots,e_k,v_k)$ along edges of
$\Proc$, with $v_k=x^{\ast}$. Each step $e_i$ is a contact operation
(\Cref{def:contact-op}) committing $\{v_{i-1},v_i\}$; the propagation
\emph{tracks} $v_0$ in that each step resolves the causal link
$\res_i\to e_{i+1}$ (information reading, \Cref{thm:three-readings}(ii)).
The \emph{committed record} of $W$ is $\rec(W)=k$.
\end{definition}

\subsection{Convergence and admissibility}

\begin{definition}[Convergent propagation]
\label{def:convergent}
A propagation $W$ is \emph{convergent to $x^{\ast}$ at tolerance $\eps$} if
its terminal state is accountable: $a(v_k,x^{\ast})\le\flo/\Omega+\eps$,
which holds with $\eps=0$ since $v_k=x^{\ast}$ gives $a=\flo/\Omega$
(the floor). More usefully, the \emph{intermediate} alignments may be
arbitrary; only the composed terminal alignment must converge.
\end{definition}

\begin{theorem}[Convergence admissibility]
\label{thm:admissibility}
A propagation $W$ is \emph{admissible} iff it is convergent to the output.
Local steps are unconstrained: an intermediate item $v_i$ may have
arbitrarily large alignment $a(v_i,x^{\ast})$ (a step ``away from'' the
output, even one corresponding to no realisable item) without rendering
$W$ inadmissible, provided the composed walk terminates at $x^{\ast}$.
Inadmissibility arises only from failure to reach $x^{\ast}$ (non-%
convergence), never from local misalignment.
\end{theorem}

\begin{proof}
A propagation is, by \Cref{def:propagation}, a walk terminating at
$x^{\ast}$; its terminal alignment is $a(x^{\ast},x^{\ast})=\flo/\Omega$
(the floor, the least possible by \Cref{lem:align-floor}), so any walk
that reaches $x^{\ast}$ is convergent ($\eps=0$). A walk that does not
reach $x^{\ast}$ has terminal item $v_k\neq x^{\ast}$, hence
$a(v_k,x^{\ast})\ge\flo/\Omega$ with strict inequality unless $v_k$ is
cut-equivalent to $x^{\ast}$; it is convergent only if it nonetheless lands
within tolerance, and inadmissible otherwise. The intermediate alignments
$a(v_i,x^{\ast})$ enter nowhere in the terminal condition, so they are
free: a step to a high-alignment (even non-realisable) $v_i$ is permitted
as long as a continuation returns the walk to $x^{\ast}$. Thus
admissibility $\iff$ convergence, with local steps unconstrained. \qed
\end{proof}

\begin{remark}[``It came from space''; the snow]
\label{rem:snow}
\Cref{thm:admissibility} is the precise sense of the following: if one
wakes to find snow and proposes that it came from space, the proposal is
not refuted by the local implausibility of the step---local alignment is
unconstrained. It is refuted only if no propagation containing that step
converges to the snow (the observed output). A locally ``impossible''
intermediate is admissible exactly when some completion through it reaches
the output; it is excluded only on non-convergence. This is the graph form
of crossing through off-target intermediates and is used in
\Cref{thm:path-opacity}.
\end{remark}

\subsection{Path opacity}

\begin{theorem}[Path opacity]
\label{thm:path-opacity}
Let $W_1,W_2$ be two propagations of the same tracked item with the same
input $v_0$ and the same output $x^{\ast}$, differing in their interior
walks. Then no invariant computed from the endpoints alone---the input
$v_0$, the output $x^{\ast}$, the terminal alignment $a(x^{\ast},x^{\ast})$,
or the minimum cut $\cut(S^{\ast}(x^{\ast}))$---distinguishes $W_1$ from
$W_2$. The interior path is free.
\end{theorem}

\begin{proof}
Every endpoint invariant is a function of $v_0$ and $x^{\ast}$ only, hence
equal on $W_1$ and $W_2$ by hypothesis. Interior data (the sequence of
intermediate vertices and committed edges) are not functions of the
endpoints---\Cref{thm:admissibility} shows the interior is unconstrained
given the endpoints---so they are not recoverable from any endpoint
invariant. Thus endpoint invariants cannot distinguish the two
propagations. \qed
\end{proof}

\begin{corollary}[The admissible set is a class, not a path]
\label{cor:gauge}
For a given $(v_0,x^{\ast},\Proc)$ the admissible propagations form an
equivalence class under ``shares endpoints and converges,'' all members
indistinguishable by endpoint invariants. The table's output is this
class, not a single chosen path.
\end{corollary}

\begin{proof}
By \Cref{thm:path-opacity} all convergent same-endpoint propagations are
endpoint-indistinguishable; by \Cref{thm:admissibility} convergence is the
sole admissibility criterion. The class of convergent same-endpoint
propagations is therefore the well-defined output. \qed
\end{proof}

\begin{figure*}[!htbp]
    \centering
    \includegraphics[width=\textwidth]{figures/panel_6_opacity.png}
    \caption{\textbf{Path Opacity: the Interior is Free.}
    (\textbf{A})~Three-dimensional rendering of two propagations sharing the same input (green) and output (orange) but distinct interior vertex orders, drawn on a circular item layout and lifted onto two path-index planes (blue and red). The two interiors differ while the endpoints coincide.
    (\textbf{B})~The endpoint invariant (the output's minimum-cut cost against the medium) across $40$ interior permutations; it is constant (points on the red dashed level), because the invariant depends only on the endpoints (Theorem~\ref{thm:path-opacity}).
    (\textbf{C})~Interior edit distance from a reference path (blue bars, left axis) growing across interior variants, against the endpoint invariant (red line, right axis) which stays flat: the interior may be edited arbitrarily without moving the endpoint invariant.
    (\textbf{D})~Distribution of the change $\Delta$ in the endpoint invariant across interior variants, a single spike at zero (red dashed): no interior rearrangement changes the endpoint invariant, so same-endpoint propagations are indistinguishable.}
    \label{fig:panel_opacity}
\end{figure*}

\subsection{Representation mobility}

\begin{definition}[Representation]
\label{def:representation}
A \emph{representation} of an item $v$ is a tuple
$\Rep(v)=(s_1,\dots,s_N)\in\R^N$ subject to the \emph{averaging
constraint}
\[
\tfrac1N\textstyle\sum_{j=1}^{N}s_j\ =\ a(v,x^{\ast}),
\]
i.e.\ the components average to the item's alignment, while individual
components are unrestricted in $\R$ (they may lie outside $(0,1]$). Two
representations of the same $v$ are \emph{equivalent}; the set of
representations of $v$ is its \emph{representation fibre}.
\end{definition}

\begin{theorem}[Representation mobility]
\label{thm:rep-mobility}
\textup{(i)} For any item $v$ and any $N\ge2$ the representation fibre is
nonempty and, for $N\ge2$, infinite: given any
$(s_1,\dots,s_{N-1})\in\R^{N-1}$ there is a unique
$s_N=N\,a(v,x^{\ast})-\sum_{j<N}s_j$ completing it. \textup{(ii)} A
propagation may change representation between consecutive steps without
changing the tracked item or the committed record: the averaging
constraint ties every representation to the same alignment, so a
representation switch commits no new cut. \textup{(iii)} Admissibility
(\Cref{thm:admissibility}) is representation-independent: it depends only on
the terminal alignment, which is the average over any representation.
\end{theorem}

\begin{proof}
(i) The averaging constraint is one linear equation in $N$ unknowns; for
$N\ge2$ its solution set is an affine subspace of dimension $N-1$, hence
infinite, and the stated $s_N$ solves it for any choice of the others.
(ii) Switching from $\Rep(v)$ to $\Rep'(v)$ (both in $v$'s fibre) leaves
$v$ and its alignment fixed (both average to $a(v,x^{\ast})$); no edge is
committed, so $\rec$ is unchanged (\Cref{def:committed}). (iii) The
terminal condition of \Cref{thm:admissibility} is $a(v_k,x^{\ast})$, which
equals the average of any terminal representation by
\Cref{def:representation}; hence admissibility is the same in every
representation. \qed
\end{proof}

\begin{corollary}[Mobility is required for heterogeneous propagation]
\label{cor:mobility-required}
A propagation whose steps are accountable only in different
representations---one contact resolved in representation $\Rep$, the next
in $\Rep'$---is admissible: by \Cref{thm:rep-mobility}(ii,iii) the
representation may be switched between the steps at no cut cost and without
affecting convergence. Hence tracking through a process that mixes contacts
of different representational type requires, and is permitted, inter-step
representation switching.
\end{corollary}

\begin{proof}
Apply \Cref{thm:rep-mobility}(ii) at the step boundary and (iii) for
admissibility. \qed
\end{proof}

\begin{figure*}[!htbp]
    \centering
    \includegraphics[width=\textwidth]{figures/panel_7_representation.png}
    \caption{\textbf{Representation Mobility and Mean Recovery.}
    (\textbf{A})~Three-dimensional scatter of $120$ three-component representations of a single item (alignment $0.5$) in $(s_1,s_2,s_3)$ space, coloured on-shell (blue, all components in $[0,1]$) or off-shell (red, some component outside): all lie on the mean-recovery plane $\tfrac13\sum s_j=0.5$, with off-shell components fully admissible (Theorem~\ref{thm:rep-mobility}).
    (\textbf{B})~Representation mean versus target alignment over $300$ random fibres and targets; every point lies on the diagonal (grey dashed): the component mean recovers the target exactly, whatever the individual components.
    (\textbf{C})~Off-shell fraction versus the component spread $M$; as the admissible spread grows the fraction of representations with an out-of-range component rises toward one, the freedom of the fibre to pass through ``impossible'' components while keeping the physical mean.
    (\textbf{D})~Record increment per operation: a representation switch (green) adds nothing to the committed record, while a cut event (red) increments it---switching representation commits no cut and leaves admissibility unchanged (Corollary~\ref{cor:mobility-required}).}
    \label{fig:panel_representation}
\end{figure*}

\subsection{The table}

\begin{definition}[Causal propagation table]
\label{def:table}
The \emph{causal propagation table} is the map
\[
\mathcal{T}\colon\ (v_0,\,I,\,\Proc)\ \longmapsto\
\acc(v_0,x^{\ast},\Proc),
\]
sending a tracked item $v_0$, input set $I\ni v_0$, and process graph
$\Proc$ with output $x^{\ast}$ to the set $\acc$ of all admissible
propagations of $v_0$ to $x^{\ast}$---the convergent same-endpoint
equivalence class (\Cref{cor:gauge}), with representations free
(\Cref{thm:rep-mobility}) and interiors free (\Cref{thm:path-opacity}).
\end{definition}

\begin{theorem}[The table is well-defined and direction-agnostic]
\label{thm:table-welldefined}
For every $(v_0,I,\Proc)$ with $x^{\ast}$ reachable from $v_0$ in $\Proc$,
$\mathcal{T}(v_0,I,\Proc)$ is a nonempty equivalence class. It is
direction-agnostic: the same class is obtained whether the process is read
forward ($v_0\to x^{\ast}$, ``synthesis'') or backward
($x^{\ast}\to v_0$, ``analysis''), because each contact operation is a
symmetric edge (\Cref{thm:symmetry}) and each step admits either reading
(\Cref{thm:reaction-measurement}).
\end{theorem}

\begin{proof}
Reachability gives at least one walk $v_0\to x^{\ast}$, which is convergent
(\Cref{thm:admissibility}), so $\acc\neq\varnothing$; by \Cref{cor:gauge}
$\acc$ is the convergent same-endpoint class. Direction-agnosticism:
reversing a walk reverses each edge, but edges are undirected
(\Cref{def:contact-op}), and each reversed step is admissible by the same
convergence criterion read from the other end (\Cref{cor:either-end}). The
forward and backward admissible sets therefore coincide. \qed
\end{proof}

\begin{definition}[Amalgamation]
\label{def:amalgamation}
The \emph{amalgamation} of $(v_0,I,\Proc)$ is the maximal set of contact
operations accounted for across $\acc$: the union, over admissible
propagations, of the resolved (information-read) contacts---``all the steps
we can account for in the propagation.'' Equivalently, it is the resolved
causal-residue chain (\Cref{cor:chain}) read as information
(\Cref{thm:three-readings}(ii)).
\end{definition}

\begin{theorem}[The amalgamation is the accounted propagation]
\label{thm:amalgamation}
The amalgamation is exactly the set of contacts whose causal links are
resolvable consistent with convergence to $x^{\ast}$. It is indifferent to
the reaction/measurement labelling of its steps
(\Cref{thm:reaction-measurement}) and to the instrument/sample labelling
of each contact's endpoints (\Cref{thm:symmetry}). Its size is the amount
of the process converted from the entropy reading to the information
reading (\Cref{thm:three-readings}); the unaccounted remainder is the
irreducible residual bounded below by the floor.
\end{theorem}

\begin{proof}
By \Cref{def:amalgamation} the amalgamation collects the resolved contacts
of admissible propagations; admissibility is convergence
(\Cref{thm:admissibility}), so a contact is in the amalgamation iff its
link is resolvable within some convergent propagation. Labelling
indifference is \Cref{thm:reaction-measurement} (reaction/measurement) and
\Cref{thm:symmetry} (instrument/sample). The size statement is
\Cref{thm:three-readings}: resolving a link moves it from entropy to
information, so the amalgamation's extent is the resolved fraction; what
cannot be resolved remains entropy, of weight $\ge\flo$ per unresolved
contact (\Cref{thm:floor}). \qed
\end{proof}

% =====================================================================
\section{The resting catalogue and the propagated table}
\label{sec:catalogue}
% =====================================================================

We now show the table's output is the same kind of object as an entry in a
catalogue of resting items---and that the familiar periodic catalogue is
the table evaluated at rest.

\subsection{The resting cut}

\begin{definition}[Resting catalogue]
\label{def:catalogue}
The \emph{resting catalogue} of a contact graph is the assignment, to each
item $v$, of its minimum cut against the medium,
$\mathrm{cell}(v):=\cut(S^{\ast}(v))$ with value $\str(v)\ge\flo$. A
catalogue \emph{cell} is the name of this cut: the invariant region by
which $v$ is told from the rest at rest (\Cref{thm:region}).
\end{definition}

\begin{theorem}[A catalogue cell is an accountable output]
\label{thm:cell-is-output}
For each item $v$, the cell $\mathrm{cell}(v)$ equals the output of the
table run on $v$ against the medium with the trivial (rest) process: i.e.\
$\mathrm{cell}(v)$ is the accountable output of the propagation that
individuates $v$ against $\med$. Equivalently, a catalogue entry is the
convergent terminal cut of the resting propagation.
\end{theorem}

\begin{proof}
Take $\Proc_{\mathrm{rest}}$ to be the contact graph with output
$x^{\ast}=v$ and the single contact $\{v,\med\}$ as available operation;
the propagation is the one cut event individuating $v$ from $\med$, whose
terminal accountable object is the minimum cut $\cut(S^{\ast}(v))$
(\Cref{def:alignment}, \Cref{thm:admissibility}). This is precisely
$\mathrm{cell}(v)$ (\Cref{def:catalogue}). \qed
\end{proof}

\begin{theorem}[The catalogue is the table at rest; the table is the
catalogue propagated]
\label{thm:catalogue-table}
The resting catalogue is the restriction of the causal propagation table
$\mathcal{T}$ to trivial (rest) processes:
$\mathrm{cell}(v)=\mathcal{T}(v,\{v\},\Proc_{\mathrm{rest}})$. Conversely,
for a nontrivial process $\Proc$, the output
$\mathcal{T}(v_0,I,\Proc)$ is an object of the same kind as a catalogue
cell---an accountable cut, the convergent terminal individuation of the
tracked item---computed for the process rather than for rest.
\end{theorem}

\begin{proof}
The forward direction is \Cref{thm:cell-is-output}. For the converse: by
\Cref{thm:admissibility} every admissible propagation terminates at the
accountable output $x^{\ast}$, whose individuating datum is the minimum cut
$\cut(S^{\ast}(x^{\ast}))$ of weight $\ge\flo$---the same type of object as
$\mathrm{cell}(\cdot)$ (\Cref{def:catalogue}), differing only in that the
cut is the terminus of a nontrivial propagation rather than of the rest
contact. Hence the table's output for a process is a catalogue-grade cut
computed under propagation. \qed
\end{proof}

\begin{corollary}[Periodic catalogue as resting table]
\label{cor:periodic}
The periodic catalogue of resting items \citep{mendeleev1869,scerri2007}
is, formally, the resting catalogue (\Cref{def:catalogue}): each entry is
the name of an item's minimum cut against the medium, itself the
accountable output of an array of resting individuations. The causal
propagation table extends it to arbitrary processes; the output of the box
is a catalogue entry of the same kind, computed for the process.
\end{corollary}

\begin{proof}
A periodic entry is a minimum-sufficient invariant of an item established
by an array of individuating operations; by \Cref{thm:cell-is-output} this
is $\mathrm{cell}(v)$, the resting table output. \Cref{thm:catalogue-table}
gives the extension. \qed
\end{proof}

\begin{remark}[Working from the catalogue, not replacing it]
\label{rem:from-catalogue}
\Cref{cor:periodic} is the precise sense in which the system \emph{works
from} the periodic catalogue rather than replacing it: the catalogue is the
fixed points of individuation (cuts at rest), and the table is the
propagation that produces fixed points. The catalogue's correctness is
inherited, not contested; the table only adds the propagated cells the
catalogue does not list because they are process-dependent.
\end{remark}

% =====================================================================
\part{The Cut}
% =====================================================================

% =====================================================================
\section{The one-use, self-blunting cut}
\label{sec:cut}
% =====================================================================

We formalise the system's emblem: the single precise cut. We show it can
be made once, that it blunts the cutter, and that the perfect repeatable
cut is the forbidden $\flo=0$ object.

\begin{definition}[The cut; the cutter]
\label{def:thecut}
A \emph{cut} is a single contact operation (\Cref{def:contact-op}): one
committed edge, individuating two items by depositing residue $\ge\flo$.
The \emph{cutter} is the endpoint read as instrument
(\Cref{thm:symmetry}); its \emph{edge-capacity} at a vertex $v$ is the set
of not-yet-committed contacts incident to $v$.
\end{definition}

\begin{theorem}[One use]
\label{thm:one-use}
A given cut---a given committed edge $e=\{u,v\}$---can be committed exactly
once. Any later cut at $u$ or $v$ is a distinct edge or the same edge at a
strictly higher committed record (\Cref{thm:monotone}), hence a different
cut event. There is no re-commitment of the same cut.
\end{theorem}

\begin{proof}
Committing $e$ increments the record (\Cref{def:committed}); a second
commitment of $e$ would either be a no-op (the edge is already committed,
so no new cut event occurs) or a distinct event at record $\rec+1$
(\Cref{thm:monotone}), which is a different cut by the record invariant.
Thus the same cut occurs once. \qed
\end{proof}

\begin{theorem}[Self-blunting]
\label{thm:blunt}
A cut deposits its residue into the cutter as much as into the cut: the
contact $\{u,v\}$ commits a boundary edge of $u$ (the cutter) no less than
of $v$ (the sample), by \Cref{thm:symmetry}. After the cut, the cutter's
edge-capacity at $u$ is strictly reduced (the committed edge is removed
from it) and the cutter's committed record has increased. Hence the cutter
is changed by the cut: it is, in the precise sense of reduced uncommitted
capacity and increased record, \emph{blunted}.
\end{theorem}

\begin{proof}
By \Cref{thm:symmetry} the operation is symmetric in $u,v$, so it commits a
boundary edge of $u$ exactly as of $v$. The committed edge leaves $u$'s
uncommitted-contact set (\Cref{def:thecut}), strictly reducing
edge-capacity, and $\rec$ increases (\Cref{thm:monotone}). Both are
irreversible (\Cref{thm:monotone}: the record cannot decrease; un-committing
is a new cut). The cutter after the cut differs from the cutter before in
capacity and record. \qed
\end{proof}

\begin{theorem}[No perfect repeatable cut]
\label{thm:no-perfect}
There is no cut that (i) separates two items at zero residue (a traceless
cut) and (ii) leaves the cutter unchanged (a repeatable instrument).
Property (i) requires $\flo=0$, forbidden by \Cref{thm:floor}; property
(ii) requires the committed record to be unchanged by a cut, forbidden by
\Cref{thm:monotone}. The perfect repeatable traceless cut is exactly the
$\flo=0$ object the floor excludes.
\end{theorem}

\begin{proof}
(i): a traceless cut deposits residue $0$, i.e.\ commits an edge of weight
$0$, contradicting $\wt(e)\ge\flo>0$ (\Cref{thm:floor}). (ii): an
unchanged cutter requires $\rec$ after $=\rec$ before, contradicting the
strict increment of \Cref{thm:monotone}. Either property alone is already
forbidden; their conjunction a fortiori. \qed
\end{proof}

\begin{remark}[Honjo Masamune: the emblem]
\label{rem:honjo}
The system's central operation is named for a blade that is ceremonial,
usable once, and blunt thereafter: the single precise cut. \Cref{thm:one-use}
is its one use; \Cref{thm:blunt} is its blunting (the residue enters the
cutter, so the instrument that cuts is spent by cutting); \Cref{thm:no-perfect}
is the reason it cannot exist as a repeatable working tool---the perfect
traceless repeatable cut is the $\flo=0$ object the whole system is built on
forbidding. What exists is not the perfect blade but the single committed,
accountable cut: cut as cleanly as the floor allows, deposit the
irreducible residue, and let that residue be both the cut's trace (the
identity) and the cause of the next cut (the propagation). Every
measurement and every reaction is such a cut (\Cref{thm:reaction-measurement}):
one-use, self-blunting, leaving the trace that is everything there is to
know. No theorem depends on the name.
\end{remark}

\begin{figure*}[!htbp]
    \centering
    \includegraphics[width=\textwidth]{figures/panel_8_convergence_cut.png}
    \caption{\textbf{Convergence Admissibility and the One-Use Cut.}
    (\textbf{A})~Three-dimensional bar chart of each item's alignment to the output $x^{*}$ (cut cost to $x^{*}$ over total weight), the output bar green and the rest blue; only the output sits at the floor, the terminal accountable state of an admissible propagation.
    (\textbf{B})~Alignment to $x^{*}$ per item with the floor alignment $\beta/\Omega$ marked (red dashed); the output (green) alone reaches the floor, so a propagation is admissible exactly when it terminates there (Theorem~\ref{thm:admissibility}).
    (\textbf{C})~Convergence outcome for three cases---reaching the output, terminating elsewhere, and a locally arbitrary interior that nonetheless converges---with admissible (blue) and expected (orange) bars coinciding in each: admissibility tracks convergence to the output and nothing else, local steps being free.
    (\textbf{D})~The one-use, self-blunting cut: the cutter's uncommitted capacity (red) decreases by one per cut while the committed record (green) increases by one, the two crossing as capacity is spent. A cut commits once and blunts the cutter; the perfect repeatable traceless cut ($\beta=0$, record unchanged) cannot occur (Theorem~\ref{thm:one-use}, Theorem~\ref{thm:blunt}, Theorem~\ref{thm:no-perfect}).}
    \label{fig:panel_convergence_cut}
\end{figure*}

% =====================================================================
\section{The master equivalence}
\label{sec:master}
% =====================================================================

We collect the system into a single biconditional, tying every structure to
the one fact $\flo>0$.

\begin{theorem}[Master equivalence]
\label{thm:master}
For a contact graph in a non-completable whole, the following are
equivalent:
\begin{enumerate}[label=\textup{(\roman*)},leftmargin=2.2em]
\item there is an identifiable item---a proper part told from the rest at
positive cost;
\item the floor is positive, $\flo>0$ (\Cref{thm:floor});
\item individuation is by negation with no selector (\Cref{thm:negation}),
identity is the invariant minimum cut and never a point
(\Cref{thm:region}), processes are monotone non-returning residue chains
(\Cref{thm:monotone,thm:propagator}), measurement and reaction are one
symmetric contact operation (\Cref{thm:reaction-measurement}), the causal
propagation table is well-defined with free representations, free interiors,
and convergence-only admissibility (\Cref{thm:table-welldefined,%
thm:rep-mobility,thm:path-opacity,thm:admissibility}), its outputs are
catalogue-grade accountable cuts (\Cref{thm:catalogue-table}), and the cut
is one-use and self-blunting with no perfect repeatable form
(\Cref{thm:one-use,thm:blunt,thm:no-perfect});
\item the perfect traceless repeatable cut does not exist
(\Cref{thm:no-perfect}).
\end{enumerate}
Each is the same fact: that telling a thing from the rest of a
non-completable whole costs a strictly positive, irreducible, one-use
residue.
\end{theorem}

\begin{proof}
$(\mathrm{i})\Rightarrow(\mathrm{ii})$: an identifiable proper part has
positive separation cost by \Cref{thm:floor-derived,thm:floor}.
$(\mathrm{ii})\Rightarrow(\mathrm{iii})$: each listed structure is proved
from $\flo>0$ in the cited theorem.
$(\mathrm{iii})\Rightarrow(\mathrm{iv})$: \Cref{thm:no-perfect} is among the
listed consequences. $(\mathrm{iv})\Rightarrow(\mathrm{i})$: if no traceless
cut exists then every separation costs $>0$, i.e.\ some item is told from
the rest at positive cost---an identifiable item. The cycle closes; the
closing sentence restates the equivalence. \qed
\end{proof}

% =====================================================================
\section{Computational validation}
\label{sec:validation}
% =====================================================================

Because every object is a finite weighted graph, every theorem is directly
checkable on explicitly constructed instances. A self-contained suite
(Python; \texttt{numpy}; \texttt{networkx} for exact max-flow minimum
cuts) implements eleven categories, one per structural result, and writes
its outcomes as JSON. All $847$ checks pass; the minimum cuts are computed
exactly (Ford--Fulkerson via \texttt{networkx}~\citep{fordfulkerson1956}),
not approximated.

\begin{table}[h]
\centering\small
\caption{Validation of the system on finite weighted graphs. Minimum cuts
are exact; all checks pass.}
\label{tab:validation}
\begin{tabular}{lll}
\toprule
\textbf{Category} & \textbf{Pass} & \textbf{Theorem checked}\\
\midrule
Floor                 & $300/300$ & \Cref{thm:floor}: $\str(v)\ge\flo$ on
                                     $300$ random contact graphs\\
Negation              & $200/200$ & \Cref{thm:negation}: complementation an
                                     involution; $U$ fixed by $\co U$\\
Identity invariant    & $101/101$ & \Cref{thm:invariant,thm:region}: cut
                                     invariant under relabelling; region-valued\\
Monotone record       & $4/4$     & \Cref{thm:monotone}: record strictly
                                     increases; no return\\
Residue propagator    & $6/6$     & \Cref{thm:propagator}: residue $>0\iff$
                                     next cut; chain halts at $0$\\
Contact symmetry      & $4/4$     & \Cref{thm:symmetry,thm:reaction-measurement}:
                                     edge symmetric; reaction $=$ measurement\\
Composition inflation & $18/18$   & \Cref{thm:trajectory-count}:
                                     $T(n,d)=d(d+1)^{n-1}$ by enumeration\\
Path opacity          & $3/3$     & \Cref{thm:path-opacity}: endpoint
                                     invariants agree across interiors\\
Representation        & $202/202$ & \Cref{thm:rep-mobility}: mean-recovery
                                     fibre; switch commits no cut\\
Convergence           & $3/3$     & \Cref{thm:admissibility}: admissible
                                     $\iff$ convergent; local steps free\\
One-use cut           & $6/6$     & \Cref{thm:one-use,thm:blunt,thm:no-perfect}:
                                     one use; blunting; no $\flo=0$ cut\\
\midrule
\textbf{Total}        & $\mathbf{847/847}$ & exact min-cut backend\\
\bottomrule
\end{tabular}
\end{table}

Three of the categories merit explicit comment, being the substantive
rather than the bookkeeping checks. \emph{(i) The floor on $300$ random
graphs.} For each of $300$ randomly weighted contact graphs (medium
adjacent to every item, random item--item contacts, all weights
$\ge\flo$), the exact minimum cut of a random item against the medium was
computed and found to satisfy $\str(v)\ge\flo$ in every case
(\Cref{thm:floor}): no sharp cut occurred. \emph{(ii) Identity is a
region.} On a graph of two strongly-internally-bound item clusters joined
to a common medium, the exact minimum cut separating one cluster from
another places a \emph{multi-item} set on one side (a five-item side in the
reported instance), confirming that the identity-bearing minimum cut is in
general not realised by a singleton (\Cref{thm:region}): identity is a
region, never a point. \emph{(iii) Composition inflation by direct
enumeration.} For every $(n,d)$ with $n\le6$, $d\in\{1,2,3\}$ small enough
to enumerate, the brute-force count of length-$n$ labelled trajectories
over $d$ move-directions with a stay-option exactly equals the closed form
$d(d+1)^{n-1}$ (\Cref{thm:trajectory-count}); larger cases inherit the
closed form. The path-opacity, representation-mobility, and convergence
categories likewise confirm, on constructed instances, that endpoint
invariants are interior-independent, that mean-recovery fibres are
nonempty with off-shell components, and that admissibility tracks
convergence to the output and nothing else.

\begin{figure*}[!htbp]
    \centering
    \includegraphics[width=\textwidth]{figures/panel_5_inflation.png}
    \caption{\textbf{Composition Inflation $T(n,d)=d(d+1)^{n-1}$.}
    (\textbf{A})~Three-dimensional surface of $\log_{10}T(n,d)$ over trajectory length $n$ and direction count $d$ (viridis); the surface rises steeply in $n$ and steepens with $d$, the exponential proliferation of distinguishable propagation trajectories with depth.
    (\textbf{B})~Trajectory count $T(n,d)$ on a logarithmic $y$-axis versus $n$ for $d=1,2,3$ (coloured lines), with directly enumerated counts (points) lying on the closed-form curves; enumeration and closed form agree exactly for every case computed.
    (\textbf{C})~Enumerated count versus closed form $d(d+1)^{n-1}$ on log--log axes; all points lie on the identity diagonal (grey dashed), the exact match of the brute-force enumeration to the formula.
    (\textbf{D})~Growth ratio $T(n,d)/T(n-1,d)$ versus $n$ for $d=1,2,3$ (coloured lines), each flat at its predicted value $d+1$ (dotted guides): the per-step branching factor is exactly $d+1$ (the $d$ moves plus the stay-option).}
    \label{fig:panel_inflation}
\end{figure*}

\begin{remark}[Scope of the suite]
\label{rem:validation-scope}
The suite verifies the theorems on instances: it confirms that the floor,
the invariance and region-valuedness of identity, the monotone record, the
residue--propagator equivalence, the contact symmetry, the inflation
count, path opacity, representation mobility, convergence admissibility,
and the one-use cut all hold as proved. It is an internal-consistency and
instance-verification check, appropriate to a system whose content is
structural; it does not, and is not intended to, supply numerical weights
for any concrete physical system, which would require a concrete weighting
$\wt$ (\Cref{sec:discussion}).
\end{remark}

% =====================================================================
\section{Discussion and scope}
\label{sec:discussion}
% =====================================================================

\subsection{What is proved, and at what level}

Every result above is a theorem about finite weighted graphs, provable from
\Cref{def:graph} and the derived floor \Cref{thm:floor}. The system is
qualitative-structural: it establishes \emph{that} individuation is a
positive-cost one-use cut, \emph{that} processes are residue chains,
\emph{that} measurement and reaction coincide, \emph{that} the table's
output is a catalogue cell, and \emph{that} the perfect cut is forbidden.
It does not compute numerical cut weights, reaction rates, or alignment
values for any concrete system; those would require a concrete weighting
$\wt$, which the framework takes as given. The relationship to quantitative
science is that of a structural theory to its instances: the framework
states what any finite system of individuated parts must satisfy and reads
off the consequences; a concrete weighting supplies the numbers.

\subsection{Honest limitations}

We mark the load-bearing soft spots. First, \emph{the non-completability
premise} (\Cref{def:noncompletable}) is an order condition we assume of the
whole; the floor is derived from it but it is itself not derived here.
Second, \emph{the alignment functional} (\Cref{def:alignment}) is one
natural cut-based scalar; other admissible scalars exist, and the
quantitative content of ``convergence'' depends on the choice, though the
qualitative theorems (admissibility $\iff$ convergence, path opacity,
representation mobility) hold for any cut-monotone alignment. Third,
\emph{representation mobility} (\Cref{thm:rep-mobility}) is proved as a
freedom (switching is permitted at no cut cost); which representation a
concrete tracking step \emph{requires} is a modelling input, not derived.
Fourth, \emph{tracking} (\Cref{thm:three-readings}) is treated as the
resolution of causal links; the framework specifies what resolution
achieves (entropy $\to$ information) but not the procedure by which a given
process is resolved, which is the empirical work the table organises rather
than performs.

\subsection{The picture in one statement}

A thing is told from the rest by a cut that costs a strictly positive,
irreducible residue. The cut can be made once; it blunts the cutter; it
leaves a trace that is the thing's identity and a residue that is the cause
of the next cut. Processes are chains of such cuts. Measurement and
reaction are the same cut, differing only in whether the trace is
accounted; instrument and sample are the two ends of the one symmetric cut.
To track an item through a process is to resolve the chain---to account for
as many cuts as converge to the output---and the accounted propagation, the
amalgamation, is the output: an accountable cut of the same kind as a
resting catalogue cell. The periodic catalogue is this table at rest; the
table is the catalogue under propagation. The perfect traceless repeatable
cut cannot exist; that it cannot is exactly why there is anything to cut,
to track, and to know.

% =====================================================================
\section{Conclusion}
\label{sec:conclusion}
% =====================================================================

We have laid out, entirely in finite weighted graphs, a complete system
generated by one derived fact---that the cut individuating a part of a
non-completable whole has strictly positive, irreducible cost $\flo>0$.
From it follow: individuation by negation; identity as the invariant
minimum cut, region-valued and never a point; the residue of each cut as
the cause of the next, making processes monotone non-returning chains; the
identity of measurement and reaction as one symmetric contact operation;
and the causal propagation table---the map from a tracked item, inputs, and
a process graph to the class of accountable propagations, with free
representations, free interiors, and convergence as the sole admissibility
criterion. The table's output is shown to be an object of the same kind as
an entry in the resting catalogue of items, so that the periodic catalogue
is the table evaluated at rest and the table is the catalogue propagated
forward. The whole is the structure of the single precise cut that can be
made once and dulls the cutter, and whose perfect traceless repeatable form
is exactly the zero-cost cut the floor forbids. The master equivalence
binds every structure to the one fact: to be told from the rest is to be
cut at irreducible cost, once.

\bibliographystyle{plainnat}
\bibliography{references}

\end{document}